% =====================================================================
% VI. Discusión crítica y limitaciones
% =====================================================================
\section{Discusión crítica y limitaciones}
\label{sec:discusion}

\textbf{La etiqueta define el techo.} Las etiquetas de Sen1Floods11 se trazaron con ayuda del
Sentinel-2 del mismo día. Parte del agua etiquetada es, por tanto, agua que el óptico ve y el SAR
no (Pakistán, Sección~\ref{sec:fallos}). Esto pone un techo al modelo solo SAR que ninguna
arquitectura puede superar, y explica a la vez la gran ganancia de la fusión. Las cifras de E6
miden ``qué tan bien se reproduce una etiqueta óptica'', no solo ``qué tan bien se detecta el
agua''.

\textbf{Lima no es un fracaso del entrenamiento, pero tampoco una prueba limpia.} La etiqueta
EMSN038 es la máxima extensión delineada con imágenes ópticas de días posteriores, no el agua
presente a la hora del paso del radar. Con un solo evento y 64 recortes con 1\,\% de agua, el IoU
es además muy sensible. Lo que sí muestran los datos multitemporales es un límite físico: el
agua urbana y de huaico se \emph{aclara}, y un modelo de una sola fecha que aprendió ``agua =
oscuro'' no puede detectarla. La salida no es más datos del mismo tipo, sino otra información:
una imagen previa para detectar cambios~\cite{bountos2024kurosiwo} o la coherencia
interferométrica~\cite{zhao2022urbanaware}.

\textbf{Tamaño de muestra y optimismo en distribución.} Entrenamos con 252 recortes y estimamos
la variabilidad con solo tres semillas. Test proviene de los mismos eventos que train, y la
autocorrelación espacial entre recortes cercanos hace esa cifra optimista. La evaluación en
Bolivia (15 recortes) es la más honesta, pero también la más ruidosa ($\pm0.026$). El IoU
agregado, a su vez, está dominado por los recortes con mucha agua y oculta la variación entre
eventos (de 0.19 a 0.89).

\textbf{Decisiones no exploradas.} Fijamos tasa de aprendizaje, profundidad, número de canales y
umbral sin búsqueda de hiperparámetros, para no ajustar nada en función de test. No probamos
ponderación de clases ni pérdida focal, aumentos radiométricos ni predicción con aumentos (TTA).
Las rotaciones de 90$^\circ$ ignoran que la geometría de observación del SAR no es invariante a
rotaciones (las sombras caen hacia el lado opuesto al sensor). En la ablación débil probamos una
sola proporción de muestreo. Por último, las corridas en Kaggle usaron estadísticas de
normalización redondeadas a tres decimales, una diferencia despreciable frente a la escala de
los datos.
